\subsection{Ball Candidate Verification by Cloth Histogram Backprojection}
\label{sec:backprojection}

This experiment investigated an independent verification stage after Hough
detection and pocket filtering. Rather than proposing more circles or changing
their geometry, it attempted to reject candidates whose pixels predominantly
resembled the table cloth. The objective was fewer false detections without
rejecting genuine balls. This is separate from the earlier cloth-removal
experiment: no pixels were removed from retained balls' recognition masks.

\subsubsection{Per-image cloth distribution and numerical scale}

Each image was converted from BGR to 8-bit HSV. A cloth-sampling mask was
constructed from the known table polygon, inset using an $11\times11$
elliptical erosion. All raw Hough circles were excluded with a four-pixel
radius margin, including candidates subsequently removed by pocket filtering.
The known pocket neighbourhoods, of radius 25 table pixels, were drawn in
table coordinates and inverse-warped to the source image for exclusion.
No annotated ball positions were used. At least 1,000 sampling pixels were
required; insufficient samples or a flat histogram caused conservative
retention of every candidate.

A two-dimensional H--S histogram used ranges $[0,180)$ and $[0,256)$.
The distribution was re-estimated for every image; no single reference-image
histogram was reused across scenes. Histogram bins were min--max normalised
to $[0,255]$, and \code{cv::calcBackProject()} produced an 8-bit similarity
map. This was converted to floating point by dividing by 255. The maps were
not separately min--max normalised afterwards. This ordering preserved
intermediate histogram scores rather than first reducing them to the
precision available from an 8-bit map in $[0,1]$. Values indicate relative
histogram similarity, not calibrated probabilities or a complete definition
of every possible cloth appearance.

\subsubsection{Candidate score and conservative retention}

Each pocket-filtered candidate retained its original centre and radius.
Its scoring mask was the original circle intersected with the image and
table mask. It was not intersected with the cloth-sampling mask, since that
mask deliberately excludes candidate interiors. Candidates with less than
80\% valid circular coverage were retained without threshold rejection.
For candidates with adequate coverage, the experiment recorded mean
similarity and the fraction $f$ of valid pixels with similarity at least
$t$. Rejection used $f>f_{\max}$, a strict comparison. Mean similarity
was diagnostic only; no edge-strength or second decision rule was added.

This rule cannot guarantee retention of a green ball whose H--S distribution
resembles cloth. The experiment therefore explicitly tracked rejected matched
balls and included green balls in the score audit. Model failures, disabled
verification and inadequate coverage had explicit retention reasons rather
than silently converting missing scores into a rejection.

\subsubsection{Recall screening and frozen evaluation}

The grid varied H bins $\{18,30\}$, S bins $\{16,32\}$,
similarity threshold $t\in\{0.6,0.7,0.8\}$ and maximum cloth fraction
$f_{\max}\in\{0.7,0.8,0.9\}$. These gave 36 configurations, evaluated
on AllBalls only. A configuration was eligible only when its model was
valid and it retained every originally matched AllBalls candidate.

All 36 configurations passed this recall screen. However, the current
AllBalls baseline had 16 matched balls and zero false detections after
pocket filtering. It therefore supplied no negative examples with which
to optimise false-positive rejection. Selection used a predeclared
conservative order: higher similarity threshold, higher maximum cloth
fraction, fewer H bins, fewer S bins, then earlier configuration ID.
Configuration 9 was frozen: 18 H bins, 16 S bins, $t=0.8$ and
$f_{\max}=0.9$. Selection was written to disk before any Table1--Table6
image or label evaluation. No alternative setting was selected after
inspecting those results. If no configuration had been eligible,
verification would have remained disabled.

The comparison included A: table masking plus raw Hough, B: A plus pocket
filtering, and C: B plus the frozen histogram verifier. Table1--Table6 were
excluded from this filter's tuning, but had already informed other project
experiments. This is a frozen-setting comparison, not pristine independent
evidence for the entire project's final generalisation accuracy.

\subsubsection{Measured result and interpretation}

% Generated by experiments/writeRecognitionReport.py from saved CSVs.
\begin{table}[htbp]
\centering
\small
\caption{Frozen cloth-verification setting on Table1--Table6. Only AllBalls was used to screen configurations; the six comparison images were not used to select a different threshold after evaluation.}
\label{tab:backprojection-results}
\begin{tabular}{lrrrrrr}
\toprule
Method & TP & FP & FN & Precision & Recall & F1 \\
\midrule
A: table mask + Hough & 86 & 2 & 4 & 97.73\% & 95.56\% & 96.63\% \\
B: A + pocket filter & 86 & 2 & 4 & 97.73\% & 95.56\% & 96.63\% \\
C: B + backprojection & 86 & 2 & 4 & 97.73\% & 95.56\% & 96.63\% \\
\bottomrule
\end{tabular}
\end{table}
\FloatBarrier

% Generated by experiments/writeRecognitionReport.py from saved CSVs.
\begin{table}[htbp]
\centering
\small
\caption{Baseline false detections under the frozen backprojection model. Fraction counts pixels with similarity at least 0.8; rejection requires more than 90\% of valid circle pixels to meet that threshold.}
\label{tab:backprojection-false-candidates}
\begin{tabular}{lrrrrc}
\toprule
Image & Candidate & Radius & Mean similarity & High-cloth fraction & Retained \\
\midrule
Table1.jpeg & 15 & 29.10 & 0.193 & 2.32\% & Yes \\
Table2.jpeg & 17 & 20.55 & 0.186 & 19.81\% & Yes \\
\bottomrule
\end{tabular}
\end{table}
\FloatBarrier


The frozen verifier rejected no candidate on any image. It preserved all
86 matched balls in Table1--Table6, but also retained both baseline false
detections. Precision remained $97.73\%$, recall $95.56\%$ and F1
$96.63\%$, with mean centre error 4.91 table pixels. AllBalls remained
16/16 with no false detections after pocket filtering. In this rerun,
raw Hough on AllBalls had one false detection, which pocket filtering
already removed. This current count should not be confused with the
historical raw-Hough tuning snapshot discussed earlier.

The two remaining false candidates were not predominantly high-similarity
cloth under the frozen model. Their high-cloth fractions were approximately
2.32\% in Table1 and 19.81\% in Table2, far below the rejection threshold
of 90\%. Thus the observed result was not a recall-versus-precision tradeoff:
the conservative rule simply did not reject either false candidate.
Detected green solid and striped balls were also retained. This observation
does not establish that arbitrary green balls are safe under another camera
or threshold configuration.

\begin{figure}[htbp]
\centering
\includegraphics[width=\linewidth]{backprojection_table1.png}
\caption{Table1 verification diagnostics: sampling exclusions, cloth-similarity
map, thresholded similarity and candidate decisions. Brightness uses the
fixed map scale, not a separate per-image display stretch. Green outlines
mean retained and red would mean rejected; these colours do not themselves
denote ground-truth correctness. The false candidate remains retained.}
\label{fig:backprojection-table1}
\end{figure}
\FloatBarrier

This first histogram-only experiment therefore does not justify adopting the
verifier. It establishes a reproducible score distribution and a recall-safe
reference setting, not an optimal rejection threshold. Additional calibration
images containing genuine post-pocket false detections would be required
to tune TP/FP separation while preserving the six images from further
threshold selection. An edge-based safeguard remains a possible separate
experiment; it was not added after looking at these evaluation outcomes.

The experiment runs with:
\begin{quote}
\code{./build-local/BallExperiments --backprojection}
\end{quote}
An optional directory argument preserves separate runs. The default directory
contains the 36-row screening grid, every AllBalls screening score, the
frozen selection, per-image and pooled A/B/C metrics, retained/rejected
candidate scores, prediction matches, rejection counts and annotation
changes. Each image has a sampling mask, histogram, YAML model, similarity
map, high-similarity map, decision overlay and comparison sheets.
Report snapshots are refreshed using:
\begin{quote}
\code{python3 experiments/writeBackprojectionReport.py}
\end{quote}
No change was made to production localisation, extraction or recognition.
